\documentclass[10pt,a4paper]{amsart}
\usepackage[margin=19mm]{geometry}
\usepackage{amsmath,amssymb,mathtools}
\usepackage[scr=rsfs,cal=euler]{mathalfa}
\usepackage{xfrac}
\usepackage[unicode,hidelinks,bookmarks=false]{hyperref}
\newcommand{\ie}{i.e.\ }
\newcommand{\cf}{cf.\ }
\newcommand{\resp}{respectively\ }
\newcommand{\Subsection}{\subsection}
\providecommand{\nobreakdash}{\nobreak\hbox{-}}
\setlength{\emergencystretch}{2em}
\multlinegap0pt
\usepackage{mathtools}
\usepackage{amssymb}
\usepackage[normalem]{ulem}
\usepackage{xcolor}
\usepackage{tikz}
\usepackage{enumerate}
\newtheorem{lemma}{Lemma}
\newcommand{\Mnat}{$\text{M}^\natural$}
\title{Tree metrics and log-concavity for matroids}
\author{Federico Ardila-Mantilla, Sergio Cristancho, Graham Denham, Christopher Eur, June Huh, Botong Wang}
\date{Source: arXiv:2601.02547v2 (CC BY 4.0)}
\begin{document}
\maketitle

\noindent\emph{Source and licence.} Tree metrics and log-concavity for matroids. Version arXiv:2601.02547v2 (CC BY 4.0), distributed by arXiv under the Creative Commons Attribution 4.0 International licence, http://creativecommons.org/licenses/by/4.0/, as recorded in the arXiv licence metadata for this exact version and shown on the abstract page https://arxiv.org/abs/2601.02547v2. Full licence notice: \texttt{licenses/arxiv-2601.02547-CC-BY-4.0.txt}.\par\smallskip

\noindent\textit{Editorial scope.} Counted unit: the article's lemma lem:ultra (source0.tex lines 823--833). The article's complete original proof of it is delivered with it (source0.tex lines 834--849, one \texttt{proof} environment of 16 lines). No mathematics is written, repaired or re-derived: the statement and the proof are the article's own lines, in the article's own notation, and no retained line is substituted or re-flowed.  No further source-local context is needed: every cross-reference of every retained line resolves inside the two delivered spans.  Nothing was omitted from any retained span, so the package carries no omission record.  The retained lines cite 1 works, whose entries are transcribed verbatim from the article's own bibliography source in the e-print and delivered with short editorial labels recorded in the item's reference mapping.  No retained line contains a figure environment, an image-inclusion command or drawing code, so no image is delivered and the package declares no asset dependency.  Editorial formatting only, in addition: an AMS \texttt{amsart} preamble that restates the article's own declarations and macros used by the retained lines, namely the environments \texttt{lemma} on separate counters, together with the standard packages those lines need.  The retained environments print in this document as Lemma 1 in order of appearance, whereas the article prints the same items with the numbering of its own full document.  The complete native source of the article as distributed in the arXiv e-print for this exact version is bundled unchanged as \texttt{sources/192162/source0.tex}.
\begin{lemma}\label{lem:ultra}
Let $\nu$ be an \Mnat-concave function on $2^E$ whose effective domain contains all subsets of $E$ with at most $2$ elements. Then,  for any $0<q\leq 1$, the function
  \[
  d(i,j)\coloneqq 
  \begin{cases}
  2q^{-\nu(i,j)+\nu(i)+\nu(j)-\nu(\varnothing)}, & \text{if $i\neq j$,}\\
  0, & \text{if $i=j$,}
  \end{cases}
  \]
  defines an ultrametric on $E$ of radius $\le 1$. %of radius $\leq 1$.  %of radius 1
\end{lemma}

\begin{proof}
  By the local exchange property for \Mnat-concave functions \cite[Theorem 6.4]{Mu03}, for any distinct elements $i,j,k$ in $E$, the maximum among the three numbers
  \[
    \nu(j,k)+\nu(i), \quad \nu(i,k)+\nu(j), \quad \nu(i,j)+\nu(k),
  \]
  is achieved at least twice. Thus, the maximum among the three numbers
  \[
    \nu(j,k) -\nu(j)-\nu(k) + \nu(\varnothing)\, \quad \nu(i,k) -\nu(i)-\nu(k) + \nu(\varnothing), \quad \nu(i,j) -\nu(i)-\nu(j) + \nu(\varnothing),
  \]
  is achieved at least twice as well. Since $0 < q \leq 1$, we get that $d$ is an ultrametric on $E$. 
  Also, the exchange property for $\nu$ gives
  \[
\nu(i,j)+\nu(\varnothing) \le \nu(i)+\nu(j) \ \ \text{for any distinct $i, j$ in $E$.}
  \]
which implies that $d$ has radius $\le 1$.
\end{proof}

\begin{thebibliography}{99}
\bibitem[Mu03]{Mu03} \bysame, \emph{Discrete {Convex} {Analysis}}, Society for Industrial and Applied Mathematics, January 2003 (en).
\end{thebibliography}
\end{document}
